\guidechapter{Making roots of unity into signed shifts}{Read before synthetic transforms and ring products.}

A transform normally multiplies values by roots of unity. In the synthetic
representation, a value is itself a short polynomial, stored as its coefficient
list. Work in
\[
 \mathcal R_r=\mathbb C[y]/(y^r+1).
\]
This means that $y^r$ is replaced by $-1$. Multiplication by $y$ shifts the
coefficient list and negates the coefficient that wraps around.

\explfigure{07-signed-shift}{A root operation becomes exact data movement.}{
For $r=4$, $y(a_0+a_1y+a_2y^2+a_3y^3)$ reduces to
$-a_3+a_0y+a_1y^2+a_2y^3$. No approximate trigonometric multiplication is
needed for this operation.}

For a power of two $t$ dividing $2r$, the element $\omega_t=y^{2r/t}$
serves as a transform root. Its powers act as signed coefficient permutations.
The source proves the needed character cancellation directly. This matters
because the quotient ring has zero divisors: dividing by a nonzero-looking
ring element is not automatically legal.

\paragraph{Keep the order the algorithm naturally produces.}
A recursive Fourier transform often produces frequencies in bit-reversed
order. For length four, that order is $0,2,1,3$. If both transformed inputs
use the same order, corresponding entries can still be multiplied. Sorting
each transform into a prettier order would add work without changing the
mathematical product. The opposite transform is written to consume the order
that is actually present.

\paragraph{Normalization is part of the interface.}
Both transforms in this manuscript contain a factor $1/t$. Consequently,
$F_t^+F_t^-=I/t$, rather than $I$. With this convention, for a length-$t$
cyclic convolution,
\[
 F_t^+\bigl((F_t^-u)\odot(F_t^-v)\bigr)=\frac{u*v}{t^2}.
\]
One factor of $t$ comes from the forward product normalization, and the
other from the normalized opposite transform. Losing either scale would
destroy exact recovery. Multidimensional versions replace $t$ by the product
of axis lengths, subject to the source's retained permutations.

\paragraph{Why polynomial multiplication does not make the proof circular.}
Pointwise entries are now polynomials, so their products still need computation.
The source packs bounded signed coefficients into larger integers with enough
separation to recover each coefficient. It handles centered extraction, signs,
and negacyclic wrapping explicitly. For these shorter integers it invokes the
already established $O(N\log N)$ multiplier, not the improved algorithm it is
trying to establish. The input length $N$ of those calls differs from the
original $n$, and all calls are charged in the final cost table.

\takeaway{A representation change can replace expensive scalar root
multiplications by exact signed shifts, but layout and scaling still cost work.}
\sourceanchor{06-transforms.tex, lem:synthetic-characters; lem:signed-ring-product}
